\section{Analýza existujúcich riešení}

Pred samotným návrhom a~implementáciou vlastného riešenia je
potrebné preskúmať, aké aplikácie a~technológie sú v~súčasnosti
dostupné v~oblasti podpory rozhodcov prostredníctvom inteligentných
zariadení. Cieľom tejto kapitoly je zmapovať existujúce riešenia
a~identifikovať ich silné a~slabé stránky.

Analýzu sme rozdelili do troch oblastí. V~prvej sa zameriavame
na rozhodcovské aplikácie pre inteligentné hodinky, ktoré sú určené
pre rôzne športy -- predovšetkým futbal, hokej či hádzanú.
V~druhej časti skúmame aplikácie špecificky zamerané na
zápasenie, no bez využitia detekcie gest. V~tretej časti
uvádzame relevantné akademické práce z~oblasti rozpoznávania
gest pomocou inerciálnych senzorov na zápästí.

\subsection{Rozhodcovské aplikácie pre inteligentné hodinky}

V~oblasti športového rozhodcovania existuje niekoľko aplikácií,
ktoré využívajú inteligentné hodinky ako nástroj na uľahčenie práce
rozhodcu počas zápasu. Tieto aplikácie sa zameriavajú najmä na
zaznamenávanie udalostí v~čase a~záznam štatistík.
Spoločným znakom všetkých nižšie uvedených riešení je, že na
ovládanie využívajú dotykový displej hodiniek -- rozhodca musí
aktívne ťukať na obrazovku, aby zaznamenal udalosť.

\subsubsection{REFSIX}
REFSIX \cite{refsix} je jedna z~najrozšírenejších aplikácií pre futbalových
rozhodcov. Je dostupná pre Apple Watch, Samsung Galaxy Watch
a~zariadenia s~Android Wear. Aplikácia umožňuje rozhodcovi
sledovať čas zápasu pomocou viacerých časomier, zaznamenávať
žlté a~červené karty spolu s~číslom hráča a~dôvodom udelenia,
a~po zápase generovať automatické správy o~priebehu stretnutia.
V~rozšírenej verzii REFSIX PRO ponúka aj sledovanie prejdenej
vzdialenosti, heatmapu pozícií rozhodcu na ihrisku a~video
analýzu. Aplikácia funguje aj bez pripojenia k~internetu --
dáta sa synchronizujú medzi hodinkami a~telefónom pred zápasom
a~po ňom.

\subsubsection{Referee Watch PRO}
Referee Watch PRO \cite{referee_watch_pro} je aplikácia pre Apple Watch určená pre
rozhodcov viacerých športov -- futbal, americký futbal, hádzanú,
rugby i hokej. Jej hlavnou funkciou je pokročilé riadenie času
zápasu s~intuitívnym ovládaním priamo z~hodiniek. Po skončení
zápasu aplikácia poskytuje analytiku výkonu vrátane heatmapy
pohybu rozhodcu po ihrisku.

\subsubsection{MatchGear}
MatchGear \cite{matchgear} je open-source aplikácia pre Android WearOS a~Apple
Watch, ktorá je určená pre rozhodcov, trénerov aj manažérov.
Aktuálne plne podporuje pozemný hokej, halový hokej, futbal,
halový futbal a~korfbal. Aplikácia umožňuje sledovanie času,
skóre, kariet a~trestných časov. Zaujímavosťou je podpora gesta
\uv{double-tap} na Apple Watch, ktoré umožňuje spustiť alebo
zastaviť čas bez nutnosti dotyku displeja. Zdrojový kód je
dostupný na platforme GitHub.

\subsubsection{Whistle}
Whistle \cite{whistle} je rozhodcovská aplikácia s~podporou inteligentných hodiniek,
telefónov a~televíznych obrazoviek. Rozhodca ovláda zápas
z~hodiniek a~všetky zmeny -- skóre, karty, prestávky -- sa
v~reálnom čase synchronizujú s~ostatnými prepojenými
zariadeniami. Aplikácia podporuje zdieľanie priebehu zápasu
prostredníctvom unikátneho kódu, vďaka čomu môžu tréneri,
asistenti rozhodcu či diváci sledovať aktuálny stav stretnutia.
Po zápase ponúka prehľad udalostí a~heatmapu pohybu rozhodcu.

\subsubsection{Spintso S1 Pro}
Na rozdiel od predchádzajúcich softvérových riešení, Spintso
S1 Pro \cite{spintso} predstavuje dedikované hardvérové zariadenie -- inteligentné
hodinky navrhnuté špeciálne pre rozhodcov. Hodinky disponujú
\Gls{gps}, meraním srdcovej frekvencie, počítadlom krokov
a~vibračnými upozorneniami pri štarte, prestávke a~konci
zápasu. Po aktivácii režimu zápasu sa dotykový displej uzamkne
a~ovládanie prebieha výhradne jedným tlačidlom, čo zabraňuje
nechceným dotykom spôsobeným napríklad dažďom alebo potom.
Hodinky je možné prepojiť s~mobilnou aplikáciou Spintso Watch,
ktorá poskytuje mapu pohybu rozhodcu a~ďalšie štatistiky.

\subsection{Aplikácie zamerané na zápasenie}

V~oblasti zápasenia existuje niekoľko aplikácií, ktoré sú však
primárne určené pre trénerov, štatistikov alebo bodových
rozhodcov sediacich pri stole -- nie pre rozhodcu na žinenke.

\subsubsection{Wrestling Scoreboard}
Wrestling Scoreboard \cite{wrestling_scoreboard} je aplikácia pre iOS, ktorá vizuálne
napodobňuje oficiálny \Gls{uww} scoreboard. Umožňuje pridávanie
a~odoberanie bodov, zaznamenávanie pasivity, napomenutí a~časomieru
s~nastaviteľnou dĺžkou polčasu. Ovládanie je čisto dotykové, a teda
rozhodca musí manuálne ťukať na tlačidlá na
obrazovke. Aplikácia nepodporuje inteligentné hodinky ani
rozpoznávanie gest.

\subsection{Akademický výskum -- rozpoznávanie gest pomocou Inertial Measurement Unit}

Popri komerčných aplikáciách prebieha aj rozsiahly akademický 
výskum zameraný na rozpoznávanie gest rúk pomocou senzorov 
pohybu umiestnených na zápästí všade tam, kde je takéto riešenie 
možné a prakticky využiteľné.

Kim a~kol. \cite{deepgesture} navrhli algoritmus rozpoznávania gest ruky
založený na hlbokom učení, konkrétne na kombinácii konvolučných
vrstiev a~rekurentných \Gls{gru} vrstiev. Vstupom boli dáta
z~gyroskopu a~akcelerometra umiestnených na zápästí. Autori
dosiahli vysokú presnosť rozpoznávania, čo potvrdilo vhodnosť
\Gls{imu} senzorov pre detekciu gest 

Xu a~kol. \cite{hand_gesture_wristband} predstavili algoritmus rozpoznávania 8~gest
ruky bežiaci na náramku s~minimálnym výpočtovým výkonom.
Navrhnutý prístup využíval kombinovanie dát z~akcelerometra
a~gyroskopu pomocou komplementárneho filtra a~definoval nový
typ príznaku nazvaný \uv{axis-crossing code}. Výsledky ukázali,
že aj na zariadeniach s~malým výkonom je možné dosiahnuť
spoľahlivé rozpoznávanie gest v~reálnom čase


V~oblasti nositeľných senzorov a~rozpoznávania ľudských aktivít
je tiež relevantný výskum z~roku 2026 \cite{wearable_sensors_ai},
ktorý navrhuje prístup s~konvolučnými neurónovými sieťami špecializovanými pre
jednotlivé typy senzorov. Výsledky jednotlivých modelov sa
následne kombinujú technikou neskorej fúzie. Tento prístup
dosahuje vyššiu presnosť ako klasické metódy, napríklad
\Gls{svm}.

\subsection{Zhrnutie a~priestor pre vlastné riešenie}

Z~vykonanej analýzy vyplýva, že existujúce rozhodcovské
aplikácie pre inteligentné hodinky sú primárne zamerané na futbal
a~ďalšie loptové športy a~ich ovládanie sa opiera o~dotykové
rozhranie -- to je v~kontexte zápasenia nepraktické, keďže
rozhodca na žinenke potrebuje mať obe ruky voľné. Aplikácie
zamerané na zápasenie zase nevyužívajú inteligentné hodinky a~slúžia
bodovým rozhodcom alebo trénerom pri stole.

Akademický výskum potvrdzuje, že rozpoznávanie gest ruky
pomocou akcelerometra a~gyroskopu na zápästí je technicky
realizovateľné s~vysokou presnosťou, žiadna z~existujúcich
prác sa však nezameriava na rozhodcovské gestá v~zápasení.
Naša práca preto vypĺňa identifikovanú medzeru a~navrhuje
riešenie, ktoré automaticky deteguje rozhodcovské gestá
prostredníctvom senzorov inteligentných hodiniek bez
nutnosti dotykového ovládania.